% ============================================================================
%  C/13-doan-bai-tom-tat — file chương
%  Ngày tạo: 2026-09-07 | Sửa gần nhất: 2026-09-07 | Ôn gần nhất: chưa
%  Dependency: C/12, B/07. Mức độ: cốt lõi.
% ============================================================================
\documentclass[../../english.tex]{subfiles}
\begin{document}
\chapter{Đoạn, bài và tóm tắt (Paragraphs, Papers, Abstracts)}
\ptrtarget{C/13}\ptrtarget{C/13-doan-bai-tom-tat}
\dependency{\ptr{C/12} cho mạch bên trong đoạn; \ptr{B/07} vì bạn đã biết cấu trúc bài báo từ phía đọc.}

\begin{tomtat}
Chương này lên hai cấp: từ câu sang đoạn, rồi sang cả bài. Ý xuyên suốt là \emph{kỳ vọng}: người đọc trong một cộng đồng khoa học mang theo một bộ khuôn rất cụ thể về chỗ nào chứa gì, và viết tốt phần lớn là đáp đúng bộ khuôn đó thay vì sáng tạo lại. Nội dung gồm: một đoạn cần gì để đứng vững; bốn câu tạo nên một bản tóm tắt; khuôn IMRAD và việc mỗi phần thật sự trả lời câu hỏi gì; ba nước đi mở bài theo mô hình CARS; và cách viết phần thảo luận mà không nói quá. Nếu chỉ nhớ một điều: khuôn không phải xiềng xích --- chúng là hợp đồng giúp người đọc tìm được thông tin, và phá chúng thì bạn trả giá chứ không phải họ.
\end{tomtat}

\begin{nentang}
\kn{câu chủ đề}{câu mở đoạn nêu ý chính của đoạn; các câu sau chống đỡ cho nó.}
\kn{IMRAD}{khuôn bài báo thực nghiệm: Introduction, Methods, Results, and Discussion.}
\kn{tóm tắt}{đoạn 150--250 từ đứng đầu bài, phải tự đứng vững một mình.}
\kn{CARS}{mô hình ba nước đi của phần mở bài: dựng bối cảnh, chỉ ra khoảng trống, chiếm chỗ trống đó.}
\kn{khoảng trống}{chỗ mà kiến thức hiện có chưa xử lý --- lý do tồn tại của bài bạn.}
\kn{đóng góp}{câu nói rõ bài này thêm gì mới; thường ở cuối phần mở bài dưới dạng danh sách.}
\end{nentang}

\begin{khoigoi}
\h{Một đoạn văn cần những gì để người đọc không lạc? Câu chủ đề có bắt buộc không?}
\h{Bản tóm tắt của bạn nên chứa những câu nào, và người đọc dùng nó để làm gì?}
\h{Vì sao gần như mọi bài báo thực nghiệm đều theo cùng một khuôn IMRAD --- và khuôn đó phục vụ ai?}
\h{Phần mở bài nên làm gì trước, và vì sao ``khoảng trống'' lại là phần khó viết nhất?}
\h{Trong phần thảo luận, làm sao nói được ý nghĩa của kết quả mà không vượt quá bằng chứng?}
\end{khoigoi}

Ở \ptr{B/07} bạn đã đọc bài báo theo ba lượt, và ba lượt đó chỉ hiệu quả vì bài báo \emph{có khuôn} --- bạn biết trước tìm cái gì ở đâu. Bây giờ đổi vai: bạn là người phải cung cấp khuôn đó.

Đây là chỗ nhiều người viết trẻ hiểu sai. Họ thấy khuôn IMRAD là hình thức bắt buộc, một thứ thủ tục phải theo, và đôi khi thấy nó gò bó sáng tạo. Cách nhìn có ích hơn: khuôn là một \emph{hợp đồng} giữa người viết và người đọc. Người đọc đồng ý bỏ công đọc bài bạn; đổi lại, bạn đồng ý đặt thông tin vào đúng chỗ họ sẽ tìm. Phá hợp đồng thì bài bạn bị đọc lướt hoặc bị hiểu sai --- và người trả giá là bạn.

\section{Đoạn: một ý, một câu mở}

Bắt đầu từ đơn vị nhỏ nhất trên câu. Câu hỏi mở đầu thứ nhất hỏi một đoạn cần gì.

Ba điều kiện, và điều kiện thứ ba hay bị bỏ quên.

\emph{Một --- đoạn có đúng một ý.} Nếu bạn không tóm được đoạn thành một câu, đoạn đó đang chứa hai ý và nên tách.

\emph{Hai --- ý đó xuất hiện sớm.} Trong văn khoa học, câu chủ đề gần như luôn là câu đầu. Đây không phải luật của tiếng Anh nói chung --- văn kể chuyện thường để ý chính ở cuối --- mà là quy ước của văn khoa học, và nó tồn tại vì người đọc \emph{đọc lướt}: họ quét câu đầu của mỗi đoạn để tìm chỗ cần đọc kỹ.

\emph{Ba --- các câu còn lại phải \emph{chống đỡ} cho ý đó, không chỉ liên quan tới nó.} Đây là điều kiện hay bị bỏ quên. Một đoạn có thể toàn câu về cùng một chủ đề mà vẫn không thành đoạn, vì không câu nào là bằng chứng, giải thích, hay ví dụ cho câu chủ đề --- chúng chỉ nằm cạnh nhau.

Vậy câu chủ đề có bắt buộc không? Trong văn khoa học: gần như có, với hai ngoại lệ. Đoạn tiếp nối mạch của đoạn trước có thể mở bằng mối nối rồi mới tới ý. Và đoạn kể một quy trình theo thứ tự thời gian có thể tự nó đã rõ. Ngoài hai trường hợp đó, nếu bạn không viết được câu chủ đề cho một đoạn, thường là vì đoạn ấy chưa có ý.

\section{Tóm tắt: bốn câu}

Câu hỏi mở đầu thứ hai. Bản tóm tắt là phần được đọc nhiều gấp hàng chục lần phần còn lại của bài, và phần lớn người đọc sẽ \emph{chỉ} đọc nó. Nó phải tự đứng vững: không trích dẫn, không viết tắt lạ, không tham chiếu tới hình trong bài.

Bốn câu, theo thứ tự này, là bộ khung dùng được cho hầu hết bài thực nghiệm; các sách hướng dẫn viết học thuật dựa trên phân tích thể loại đưa ra những bộ khung rất gần\cite{swalesfeak2012}.

\begin{center}\footnotesize
\rowcolors{2}{white}{tblalt}
\begin{longtable}{@{}L{1.4cm}L{4.3cm}L{8.5cm}@{}}
\toprule
\textbf{Câu} & \textbf{Trả lời câu hỏi} & \textbf{Mẫu}\\
\midrule
\endfirsthead
\toprule
\textbf{Câu} & \textbf{Trả lời câu hỏi} & \textbf{Mẫu (tiếp)}\\
\midrule
\endhead
1 & Bài toán là gì, vì sao đáng quan tâm & \emph{X is essential for Y, but current methods fail when \ldots}\\
2 & Khoảng trống cụ thể & \emph{Existing approaches assume Z, which does not hold in \ldots}\\
3 & Bạn làm gì & \emph{We propose \ldots, which \ldots}\\
4 & Kết quả, có số & \emph{On \ldots, our method reduces error by 32\% while running at 30 Hz.}\\
\bottomrule
\end{longtable}
\end{center}

Hai lỗi phổ biến. Lỗi thứ nhất là dành ba câu đầu cho bối cảnh chung chung và chỉ còn một câu cho đóng góp --- người đọc bỏ đi trước khi tới chỗ quan trọng. Lỗi thứ hai là câu 4 không có số: \emph{results show the effectiveness of the proposed method} không nói gì cả, và người đọc có kinh nghiệm đọc nó như một tín hiệu rằng kết quả yếu (\ptr{B/08}).

\begin{capdoi}[Câu kết quả]
\truoc{Experimental results demonstrate the effectiveness of the proposed approach.}
\sau{On the EuRoC sequences, our method reduces absolute trajectory error from 3.4 cm to 2.3 cm while keeping the same runtime.}
\vico{Câu trên không thể sai nên cũng không mang thông tin. Câu dưới nêu tập dữ liệu, chỉ số, mức cải thiện và cái không đánh đổi.}
\end{capdoi}

\section{IMRAD: mỗi phần trả lời một câu hỏi}

Câu hỏi mở đầu thứ ba hỏi vì sao khuôn này phổ biến đến vậy và nó phục vụ ai. Câu trả lời ngắn: nó phục vụ \emph{người đọc chọn lọc}, tức là gần như mọi người đọc. Bốn phần tương ứng bốn câu hỏi mà một người đánh giá công trình bắt buộc phải hỏi, theo đúng thứ tự đó.

\begin{center}\footnotesize
\rowcolors{2}{white}{tblalt}
\begin{longtable}{@{}L{2.4cm}L{5.0cm}L{7.0cm}@{}}
\toprule
\textbf{Phần} & \textbf{Trả lời} & \textbf{Lỗi thường gặp}\\
\midrule
\endfirsthead
\toprule
\textbf{Phần} & \textbf{Trả lời} & \textbf{Lỗi thường gặp (tiếp)}\\
\midrule
\endhead
Mở bài & Vì sao phải làm việc này? & Bối cảnh quá rộng; không nêu rõ khoảng trống\\
Phương pháp & Bạn làm chính xác cái gì? & Thiếu chi tiết để lặp lại; trộn lẫn kết quả vào\\
Kết quả & Quan sát được gì? & Diễn giải sớm; chỉ đưa số có lợi\\
Thảo luận & Nó có nghĩa gì, và giới hạn ở đâu? & Nói quá; giấu hạn chế; lặp lại phần kết quả\\
\bottomrule
\end{longtable}
\end{center}

Ranh giới quan trọng nhất --- và bị vi phạm nhiều nhất --- là giữa Kết quả và Thảo luận. Kết quả nói \emph{cái đo được}; Thảo luận nói \emph{nó nghĩa là gì}. Trộn hai thứ khiến người đọc không tách được quan sát khỏi diễn giải, và đó chính là kỹ năng mà \ptr{B/06} dạy họ phải làm. Nếu bạn làm khó việc đó, người đọc cẩn thận sẽ mất niềm tin.

\begin{cannho}
Có một phép thử nhanh cho phần Phương pháp: đưa cho một đồng nghiệp cùng ngành và hỏi ``bạn có dựng lại được thí nghiệm này không?''. Nếu họ phải hỏi lại quá ba câu, phần đó chưa đủ. Đây là tiêu chí khách quan thay cho cảm giác ``đã viết đủ chưa''.
\end{cannho}

\section{Mở bài: ba nước đi}

Câu hỏi mở đầu thứ tư. Phần mở bài là phần khó nhất với hầu hết người viết, và có một mô hình mô tả rất chính xác cái mà các bài báo thực tế đang làm --- mô hình CARS của Swales\cite{swales1990}. Nó không phải một khuôn được nghĩ ra rồi áp xuống; nó là kết quả của việc \emph{quan sát} hàng loạt bài báo và rút ra cấu trúc chung. Đó là lý do nó dùng được: bạn đang đáp lại kỳ vọng có thật.

\begin{figure}[H]\centering
\begin{tikzpicture}[node distance=0.5cm,
  m/.style={draw=steel,fill=bluebg,rounded corners=2pt,inner sep=5pt,align=left,font=\small,text width=10.4cm},
  ar/.style={-{Latex[length=2.2mm]},draw=accent,thick},
  lb/.style={font=\scriptsize\itshape,color=tiercol}]
\node[m] (m1) {\textbf{Nước đi 1 --- Dựng bối cảnh.} Lĩnh vực này quan trọng; đây là những gì đã biết. \emph{Ngắn:} 1--2 đoạn, không phải một bài tổng quan.};
\node[m,below=of m1] (m2) {\textbf{Nước đi 2 --- Chỉ ra khoảng trống.} Nhưng phương pháp hiện có giả định \(Z\), và giả định đó hỏng khi \ldots \emph{Đây là phần khó nhất.}};
\node[m,below=of m2] (m3) {\textbf{Nước đi 3 --- Chiếm chỗ trống.} Bài này làm \ldots\ Đóng góp gồm: (i) \ldots\ (ii) \ldots\ (iii) \ldots};
\draw[ar] (m1) -- (m2);
\draw[ar] (m2) -- (m3);
\node[lb,right=0.2cm of m2.east,text width=2.2cm,align=left] {khoảng trống quyết định giá trị bài};
\end{tikzpicture}
\caption{Ba nước đi của phần mở bài. Nước đi 2 là bản lề: nó vừa biện minh cho bài, vừa định nghĩa trước tiêu chí mà người đọc sẽ dùng để đánh giá phần kết quả.}
\label{fig:cars-en}
\end{figure}

Vì sao nước đi 2 khó nhất? Vì nó đòi hai việc cùng lúc mà nhiều người mới chỉ làm được một. Bạn phải \emph{biết đủ} tài liệu để nói được cái gì chưa có --- đó là việc đọc. Và bạn phải nói điều đó \emph{mà không hạ thấp} công trình của người khác --- đó là việc xã giao. Người viết thiếu kinh nghiệm hoặc né hẳn (viết bối cảnh rồi nhảy thẳng sang đóng góp, khiến bài không có lý do tồn tại), hoặc nói quá gay gắt (\emph{previous work has completely ignored\ldots}, gần như luôn sai và gây khó chịu cho chính những người sẽ phản biện bài của bạn).

\begin{mauau}[Nêu khoảng trống một cách lịch sự]
\mc{Giới hạn điều kiện}{These methods work well when X holds, but X rarely holds in \ldots}{Công nhận trước, giới hạn sau}
\mc{Chưa được xét}{The effect of X on Y has received less attention.}{Không nói ai sai, chỉ nói còn thiếu}
\mc{Đánh đổi chưa giải}{Existing approaches trade X for Y; we ask whether both can be kept.}{Đặt bài của bạn thành câu hỏi}
\mc{Bằng chứng còn mỏng}{The evidence for X comes mainly from \ldots, leaving open whether \ldots}{Mạnh mà vẫn công bằng}
\mc{Mở rộng chứ không bác bỏ}{We build on \textup{[7]} and extend it to the case where \ldots}{An toàn nhất khi bài bạn là cải tiến}
\end{mauau}

\section{Thảo luận: nói ý nghĩa mà không nói quá}

Câu hỏi mở đầu thứ năm. Phần thảo luận là chỗ dễ mất uy tín nhất, vì đây là lúc bạn rời khỏi số liệu.

Bốn việc mà một phần thảo luận tốt làm, theo thứ tự\cite{alley2018}. \emph{Một:} nêu lại kết quả chính bằng một câu, không lặp bảng số. \emph{Hai:} giải thích \emph{vì sao} nó xảy ra --- đây là phần có giá trị nhất và cũng là phần dễ đoán mò nhất, nên mức cam kết phải khớp bằng chứng (\ptr{C/14}). \emph{Ba:} nêu hạn chế \emph{trước khi} người phản biện nêu hộ bạn. \emph{Bốn:} nói kết quả này mở ra cái gì.

Về việc ba, có một hiểu lầm đáng gỡ: nhiều người tránh nêu hạn chế vì sợ bài yếu đi. Thực tế ngược lại. Một bài tự nêu ranh giới áp dụng cho thấy tác giả hiểu công trình của mình, và nó ngăn người đọc tự tưởng tượng ra hạn chế còn tệ hơn. Điều cần tránh không phải là nêu hạn chế mà là nêu theo kiểu tự vệ hoặc giấu chúng ở câu cuối cùng.

\section{Gốc rễ: khuôn từ đâu ra}

Bài báo khoa học không phải lúc nào cũng có khuôn IMRAD. Các bài của thế kỷ 17 và 18 gần với thư từ và tường thuật hơn: tác giả kể lại họ đã làm gì theo trình tự thời gian, xen lẫn suy nghĩ, và kết luận nằm rải rác. Cấu trúc bốn phần chỉ trở nên phổ biến trong thế kỷ 20, và trở thành gần như bắt buộc trong y sinh và các ngành thực nghiệm vào nửa sau thế kỷ đó.

Động lực rất thực dụng và có ba mặt. Thứ nhất là \emph{khối lượng}: khi số bài xuất bản tăng nhanh, không ai đọc hết được, nên người đọc cần một cách quét nhanh --- và điều đó chỉ khả thi nếu mọi bài đặt cùng loại thông tin ở cùng chỗ. Thứ hai là \emph{phản biện}: người phản biện phải kiểm tra được ba việc tách bạch --- phương pháp có hợp lý không, số liệu có ủng hộ kết luận không, kết luận có vượt quá số liệu không --- và họ chỉ làm được nếu ba thứ đó nằm ở ba chỗ khác nhau. Thứ ba là \emph{tính lặp lại}: tách phần phương pháp ra riêng khiến việc thiếu chi tiết trở nên lộ liễu.

Mô hình CARS ra đời theo một logic hoàn toàn khác, và điều này đáng chú ý. Swales không thiết kế khuôn; ông \emph{phân tích} các phần mở bài có thật và tìm ra cấu trúc lặp lại. Nói cách khác, CARS là một mô tả thực nghiệm về hành vi cộng đồng, không phải một quy định. Đó cũng là lý do nó thuyết phục hơn nhiều so với các danh sách ``bảy bước viết mở bài'' xuất hiện trong sách hướng dẫn.

Một hệ quả cần nhớ: vì khuôn là quy ước cộng đồng, chúng \emph{khác nhau giữa các ngành}. Toán học thuần tuý không dùng IMRAD. Nhiều hội nghị máy tính đặt phần công trình liên quan sau phần phương pháp thay vì trước. Cách đúng để biết khuôn của cộng đồng bạn không phải đọc sách hướng dẫn chung\cite{swalesfeak2012,booth2016} mà là đọc năm bài gần đây ở đúng nơi bạn định gửi và ghi lại cấu trúc của chúng.

\begin{gocre}
\gr{Thế kỷ 17--18 --- bài báo dạng tường thuật}{Báo cáo cho một cộng đồng nhỏ, ai cũng biết nhau.}{Kể theo trình tự thời gian; không có khuôn cố định.}
\gr{Thế kỷ 20 --- IMRAD lan rộng}{Số bài tăng nhanh, không ai đọc hết.}{Cùng loại thông tin ở cùng chỗ \(\Rightarrow\) quét nhanh được.}
\gr{Chuẩn hoá phản biện}{Người phản biện cần kiểm ba việc tách bạch.}{Phương pháp, kết quả, thảo luận tách hẳn; nói quá trở nên dễ phát hiện.}
\gr{Yêu cầu lặp lại thí nghiệm}{Kết quả không lặp lại được thì vô giá trị.}{Phần phương pháp thành mục riêng, thiếu chi tiết là lộ ngay.}
\gr{1990 --- Swales và CARS\cite{swales1990}}{Các phần mở bài thật sự làm gì?}{Mô tả thực nghiệm ba nước đi --- không phải quy định, nên mới đáng tin.}
\gr{Khác biệt giữa ngành}{Một khuôn có hợp mọi ngành không?}{Không: toán không dùng IMRAD, các hội nghị đặt phần liên quan ở chỗ khác \(\Rightarrow\) học khuôn từ chính nơi bạn gửi bài.}
\end{gocre}

\section{Liên hệ ngang}

\begin{lienhe}
\lh{Lập trình}{Quy ước dự án và bố cục thư mục}{Cùng lý do tồn tại: người mới tìm được thứ mình cần vì mọi dự án đặt nó ở cùng chỗ. Phá quy ước thì chính bạn trả giá bằng thời gian của người đọc mã.}
\lh{Báo chí}{Kim tự tháp ngược}{Cùng nguyên tắc với tóm tắt: thông tin quan trọng nhất lên đầu vì người đọc có thể dừng bất cứ lúc nào. Khác với văn kể chuyện, nơi kết luận để dành cuối.}
\lh{Kiến trúc phần mềm}{Tài liệu quyết định thiết kế}{Cũng có khuôn ba phần: bối cảnh, khoảng trống, quyết định --- gần như CARS. Cùng lý do: người đọc sau này cần biết \emph{vì sao}, không chỉ \emph{cái gì}.}
\lh{Pháp lý}{Bản án và bố cục lập luận}{Tách bạch sự kiện, luật áp dụng, và kết luận --- giống việc tách Kết quả khỏi Thảo luận, và cũng vì lý do kiểm tra được.}
\lh{Trình bày miệng}{Slide đầu tiên của báo cáo}{Bốn câu tóm tắt ở mục 2 dùng nguyên được cho slide mở đầu; khán giả cũng quyết định ở lại hay không trong 30 giây đầu.}
\lh{Nhân học}{Quy ước của một cộng đồng}{Khuôn viết là một tập tục nhóm; hiểu như vậy giải thích vì sao chúng khác nhau giữa ngành và vì sao phá chúng bị phạt về mặt xã hội chứ không phải logic.}
\end{lienhe}

\begin{traloi}
\tl{Ba điều kiện: đoạn có đúng một ý (nếu không tóm được thành một câu thì nên tách); ý đó xuất hiện sớm, thường ở câu đầu; và các câu còn lại \emph{chống đỡ} cho ý đó chứ không chỉ liên quan tới nó --- điều kiện thứ ba hay bị bỏ quên. Câu chủ đề gần như bắt buộc trong văn khoa học vì người đọc quét câu đầu mỗi đoạn; hai ngoại lệ là đoạn tiếp nối mạch đoạn trước và đoạn kể quy trình theo thời gian.}
\tl{Bốn câu: bài toán và vì sao đáng quan tâm; khoảng trống cụ thể; bạn làm gì; kết quả kèm số. Người đọc dùng nó để quyết định có đọc tiếp không, và phần lớn sẽ chỉ đọc nó --- nên nó phải tự đứng vững, không trích dẫn, không viết tắt lạ, không tham chiếu hình. Hai lỗi hay gặp: dồn ba câu cho bối cảnh chung chung, và câu kết quả không có số.}
\tl{Nó phục vụ người đọc chọn lọc, tức gần như mọi người đọc. Bốn phần ứng với bốn câu hỏi bắt buộc khi đánh giá một công trình: vì sao làm, làm cái gì, quan sát được gì, nó nghĩa là gì. Ba động lực lịch sử: khối lượng bài tăng nên cần quét nhanh; phản biện cần kiểm ba việc tách bạch; và yêu cầu lặp lại thí nghiệm khiến phần phương pháp phải đứng riêng.}
\tl{Ba nước đi theo mô hình CARS: dựng bối cảnh (ngắn, 1--2 đoạn), chỉ ra khoảng trống, rồi chiếm chỗ trống đó bằng đóng góp của bạn. Nước đi 2 khó nhất vì đòi hai việc cùng lúc: biết đủ tài liệu để nói được cái gì chưa có, và nói điều đó mà không hạ thấp công trình người khác --- những người rất có thể sẽ phản biện bài bạn. Né hẳn thì bài không có lý do tồn tại; nói gay gắt thì gần như luôn sai.}
\tl{Bốn việc theo thứ tự: nêu lại kết quả chính bằng một câu; giải thích vì sao nó xảy ra, với mức cam kết khớp bằng chứng (\ptr{C/14}); nêu hạn chế trước khi người phản biện nêu hộ; và nói kết quả mở ra cái gì. Nêu hạn chế không làm bài yếu đi mà cho thấy bạn hiểu công trình của mình, đồng thời ngăn người đọc tự tưởng tượng ra hạn chế tệ hơn.}
\end{traloi}

\begin{dongian}
Khi viết một bài báo hay một báo cáo, có một điều nên hiểu trước mọi mẹo viết: người đọc đến với bài của bạn kèm sẵn một danh sách chỗ nào chứa gì. Họ biết phần đầu nói vì sao đề tài đáng làm, phần giữa nói bạn làm gì và đo được gì, phần cuối nói nó có nghĩa gì. Họ không đọc từ đầu tới cuối --- họ nhảy tới chỗ họ cần.

Nên viết tốt phần lớn là \emph{đặt đúng thứ vào đúng chỗ}. Nghe không hấp dẫn, nhưng nó tiết kiệm cho bạn rất nhiều công.

Với một đoạn văn: một đoạn nên có đúng một ý, và ý đó nên nằm ở câu đầu tiên. Lý do rất thực tế: người đọc lướt qua câu đầu của từng đoạn để tìm chỗ đáng dừng lại. Nếu bạn giấu ý ở câu cuối, họ bỏ qua đoạn đó. Và nếu bạn không viết nổi câu đầu tóm ý cả đoạn, thường là đoạn đó đang chứa hai ý --- hãy tách ra.

Với bản tóm tắt: bốn câu là đủ. Vấn đề là gì. Cái gì còn thiếu. Bạn làm gì. Kết quả ra sao --- và câu cuối này bắt buộc phải có số. \emph{Results show the effectiveness of our method} là câu không thể sai, nên cũng không nói gì; người đọc có kinh nghiệm đọc nó như dấu hiệu kết quả yếu.

Với phần mở bài: ba bước. Đây là chuyện gì và người ta đã biết gì (ngắn thôi). Còn thiếu chỗ nào. Bài này lấp chỗ đó ra sao. Bước giữa là bước khó nhất, vì bạn phải nói người khác chưa làm gì mà không được nói họ dở --- rất có thể chính họ sẽ chấm bài bạn. Cách an toàn: nói về \emph{điều kiện}, không nói về người. ``Các phương pháp này hoạt động tốt khi X, nhưng X hiếm khi đúng trong trường hợp\ldots''

Và một điều nhiều người làm ngược: hãy tự nêu hạn chế của công trình mình. Không phải để khiêm tốn. Mà vì nó cho thấy bạn hiểu việc mình làm, và nó ngăn người đọc tự nghĩ ra hạn chế còn tệ hơn cái thật.
\motcau{Khuôn không gò bạn --- nó là chỗ người đọc sẽ tìm.}
\chuadongian{Khuôn khác nhau giữa các ngành, nên không có bộ khuôn chung nào thay được việc đọc năm bài gần đây ở đúng nơi bạn định gửi.}
\end{dongian}

\begin{onbai}
\ar[3]{Ba điều kiện của một đoạn}{Đúng một ý; ý ở câu đầu; các câu sau chống đỡ chứ không chỉ liên quan.}
\ar[3]{Bốn câu tóm tắt}{Bài toán \(\rightarrow\) khoảng trống \(\rightarrow\) bạn làm gì \(\rightarrow\) kết quả có số.}
\ar[3]{Câu kết quả phải có số}{\emph{Demonstrate the effectiveness} là câu không thể sai nên không mang tin.}
\ar[3]{Bốn phần IMRAD}{Vì sao làm; làm cái gì; quan sát được gì; nó nghĩa là gì.}
\ar[3]{Ranh giới Kết quả--Thảo luận}{Kết quả nói cái đo được; Thảo luận nói nó nghĩa gì. Trộn lẫn làm mất niềm tin.}
\ar[3]{Ba nước đi CARS}{Dựng bối cảnh \(\rightarrow\) chỉ ra khoảng trống \(\rightarrow\) chiếm chỗ trống.}
\ar[2]{Vì sao nước đi 2 khó}{Đòi vừa biết đủ tài liệu vừa nói được mà không hạ thấp người khác.}
\ar[2]{Cách nêu khoảng trống an toàn}{Nói về điều kiện, không nói về người: ``works well when X, but X rarely holds\ldots''}
\ar[2]{Bốn việc của phần thảo luận}{Nêu lại kết quả chính; giải thích vì sao; nêu hạn chế; nói nó mở ra gì.}
\ar[2]{Vì sao nên tự nêu hạn chế}{Cho thấy bạn hiểu công trình mình và ngăn người đọc tưởng tượng hạn chế tệ hơn.}
\ar[2]{Phép thử phần phương pháp}{Đồng nghiệp cùng ngành dựng lại được thí nghiệm mà hỏi lại không quá ba câu.}
\ar[1]{Khuôn là quy ước cộng đồng}{Khác nhau giữa ngành; học bằng cách đọc năm bài gần đây ở đúng nơi bạn gửi.}
\end{onbai}

\begin{hoidap}
\qa{Tôi thấy khuôn IMRAD gò bó. Có được phá không?}{Được, nhưng hãy biết cái giá. Người đọc quét bài bạn theo khuôn quen; nếu bạn đặt thông tin ở chỗ khác, họ không tìm thấy và kết luận là bài thiếu. Trong thực tế, những bài phá khuôn thành công thường là bài của tác giả đã có uy tín, hoặc là bài thuộc thể loại khác hẳn (bài quan điểm, bài tổng quan). Với một bài thực nghiệm thông thường, chi phí gần như luôn lớn hơn lợi ích.}
\qa{Phần công trình liên quan nên đặt ở đâu?}{Tuỳ cộng đồng, và đây là ví dụ rõ nhất cho việc khuôn khác nhau giữa ngành. Đặt sau phần mở bài là truyền thống và tiện khi bạn cần nhiều nền để hiểu phương pháp. Đặt sau phần phương pháp có lợi khi bạn muốn người đọc hiểu ý tưởng của bạn trước rồi mới so sánh --- nhiều hội nghị máy tính theo cách này. Cách quyết: xem năm bài gần đây ở đúng nơi bạn gửi.}
\qa{Đoạn của tôi nên dài bao nhiêu?}{Không có con số đúng, nhưng có hai dấu hiệu. Đoạn dài quá nửa cột báo thường chứa hai ý --- thử tóm nó thành một câu, nếu phải dùng chữ ``và'' thì tách. Đoạn chỉ một câu thường là dấu hiệu ý đó chưa được chống đỡ, trừ khi bạn cố ý dùng nó để nhấn. Điều đáng lo hơn độ dài là điều kiện thứ ba ở mục 1: các câu có thật sự chống đỡ câu chủ đề không.}
\qa{Nên viết phần nào trước?}{Phần lớn người viết có kinh nghiệm không viết theo thứ tự bài. Trình tự thường dùng: Phương pháp và Kết quả trước (vì chúng đã cố định), rồi Thảo luận, rồi Mở bài, và Tóm tắt viết \emph{cuối cùng}. Lý do: phần mở bài phải nêu khoảng trống khớp với cái bài thật sự làm, và bạn chỉ biết chắc điều đó sau khi đã viết xong kết quả. Viết mở bài trước hay dẫn tới việc hứa một đằng làm một nẻo.}
\qa{Tôi có nên dùng tiêu đề phụ trong phần thảo luận không?}{Nếu phần đó dài hơn khoảng một trang thì nên. Tiêu đề phụ giúp người đọc quét, và quan trọng hơn, chúng buộc bạn phải tách bạch bốn việc ở mục 5 --- nhất là tách phần hạn chế ra thành một mục riêng thay vì giấu nó trong một câu. Nhiều cộng đồng có quy ước riêng về việc này, nên vẫn áp dụng nguyên tắc xem bài gần đây.}
\qa{Câu chủ đề có phải luôn là câu tuyên bố khô khan không?}{Không. Nó có thể là một câu hỏi mà đoạn sẽ trả lời, hoặc một tuyên bố có mối nối ở đầu theo nguyên tắc \ptr{C/12}. Yêu cầu duy nhất là sau khi đọc câu đó, người đọc biết đoạn này định làm gì. Một câu chủ đề tốt vừa nối với đoạn trước ở đầu vừa nêu ý mới ở cuối --- đúng cấu trúc cũ trước mới sau.}
\qa{Viết tóm tắt trước khi viết bài có phải ý hay không?}{Viết một bản \emph{nháp} tóm tắt từ sớm thì rất hữu ích, nhưng vì lý do khác với bạn nghĩ: nó buộc bạn phát biểu khoảng trống và đóng góp thành câu, và thường làm lộ ra rằng bạn chưa rõ chúng. Hãy coi bản nháp đó như một công cụ suy nghĩ, và viết lại tóm tắt thật sau khi có kết quả cuối cùng --- vì câu thứ tư cần số thật.}
\end{hoidap}

\begin{baitap}
\bt[1]{Lấy 5 đoạn trong bài của bạn và viết câu tóm mỗi đoạn thành một câu}{Bài viết của bạn}{Đoạn nào phải dùng chữ ``và'' để tóm thì tách đôi.}
\bt[1]{Phân tích 3 bản tóm tắt theo bốn câu ở mục 2}{Bài báo ngành}{Đánh dấu câu nào thiếu; xem câu kết quả có số không.}
\bt[2]{Viết lại bản tóm tắt của bạn đúng bốn câu, câu cuối có số}{Bài viết của bạn}{Giới hạn 200 từ; không trích dẫn, không viết tắt lạ.}
\bt[2]{Phân tích 3 phần mở bài theo ba nước đi CARS}{Bài báo ngành}{Chỉ rõ câu nào là nước đi 2; bài nào thiếu nước đi 2 thì ghi lại.}
\bt[2]{Viết 5 cách nêu khoảng trống cho đề tài của bạn, dùng mẫu ở mục 4}{Bài viết của bạn}{Không câu nào được nói người khác dở.}
\bt[3]{Lấy một bài báo, tách các câu trong phần Thảo luận thành ``quan sát'' và ``diễn giải''}{Bài báo ngành}{Câu nào nằm giữa hai loại thì đó là chỗ tác giả nói quá.}
\bt[3]{Viết mục hạn chế cho công trình của bạn, ít nhất ba hạn chế thật}{Bài viết của bạn}{Mỗi hạn chế kèm một câu nói nó ảnh hưởng kết luận nào.}
\end{baitap}

\section*{Kết chương và đọc tiếp}

Ba chương vừa rồi đi từ câu lên đoạn lên bài, và mỗi cấp có một nguyên tắc riêng: ai làm gì, cũ trước mới sau, và đáp đúng kỳ vọng về khuôn.

Nhưng có một chiều cắt ngang cả ba cấp mà ta mới chỉ chạm tới. Ở mục 5 vừa rồi, việc ``giải thích vì sao kết quả xảy ra'' đòi bạn nói một điều mà bằng chứng chỉ ủng hộ một phần --- và cách bạn diễn đạt mức tin đó quyết định người đọc có tin bạn hay không. Đây là chỗ người viết không bản ngữ hay sai theo cả hai chiều: mạnh quá thành thiếu cơ sở, nhẹ quá thành không khẳng định gì. Chương \ptr{C/14} là mặt viết của \ptr{B/08}.
\ifSubfilesClassLoaded{\printbibliography[title={Nguồn}]}{}
\end{document}
